% =====================================================================
%  Section 2: Related Work
% =====================================================================
\section{Related Work}
\label{sec:related}

Instruction-tuning data selection has been studied through both
static and dynamic criteria. Quality-based methods assign each
candidate a global utility score, either using a strong external
judge, as in AlpaGasus~\cite{alpagasus}, or a learned quality
estimator, as in CPQS~\cite{cpqs2026}. These methods are simple and
task-agnostic, but a single pool-level ranking cannot adapt to a
target capability or to the model's changing state during training.

Target-conditioned methods define usefulness relative to examples
that represent a desired capability. Influence- and gradient-based
approaches such as LESS~\cite{less} score candidates by their
gradient similarity to a small validation or few-shot target set,
while Approx-LESS~\cite{approxless} improves scalability by
predicting influence scores from a subset of gradient features.
Another line uses representation geometry rather than gradients:
NAIT~\cite{nait2026} extracts hidden-state capability directions
from in-domain seed examples and selects candidates aligned with
those directions. These approaches can produce capability-aligned
subsets, but their selection criterion is fixed before training and
depends on seed or validation examples supplied by the practitioner.
\paragraph{Dynamic data selection for RL fine-tuning.}
A closely related but experimentally non-comparable line studies
\emph{dynamic} data selection for RL fine-tuning with verifiable
rewards (RLVR). DOTS~\cite{sun2025dots} targets GRPO training by
selecting questions whose \emph{adaptive difficulty}---the
policy-dependent rollout failure rate---is close to one half, where
binary reward variation yields the strongest group-relative gradient
signal. This makes DOTS an effective dynamic question-selection
method for RLVR, but its signal is defined through rollout groups,
verifiable rewards, and a learned difficulty predictor.

DOTS and \textsc{TADS} both belong to dynamic data selection, but
they are not interchangeable baselines. DOTS selects RLVR questions
under GRPO; \textsc{TADS} selects supervised instruction--response
pairs under SFT. Applying DOTS to our setting would require
converting SFT data into verifiable RLVR questions and adding reward
design, rollout generation, and difficulty prediction, thereby
changing the objective, data unit, supervision source, and compute
profile. We therefore treat DOTS as the closest dynamic selector in
a different training regime, not as a directly runnable baseline for
dynamic SFT data selection.

\paragraph{Positioning of \textsc{TADS}.}
\textsc{TADS} connects static representation-based selection and
dynamic data selection by making dynamic instruction-tuning
selection representation-aware. Static representation-based methods
use hidden-state geometry, but their capability directions are fixed
before training and typically depend on seed or target examples.
Dynamic selectors adapt to the current model state, but their
signals are usually scalar utilities---loss, uncertainty, reward, or
policy weights---rather than structural directions in the model's
evolving representation space.

Our implementation uses an advantage-estimator dynamic selection pipeline
that supplies two time-dependent signals for each candidate: a
composite reward $\mathcal{R}_i^{(t)}$ and a learned advantage estimate
$a_i^{(t)}$. We combine them as a fixed
\emph{reward--advantage diagnostic consensus} ranking,
$\mathcal{R}_i^{(t)}a_i^{(t)}$: a candidate receives a high base
rank only when both the scalar reward signal and the advantage estimate
support its selection. This is a selection-time ranking rule, not an
additional training objective. \textsc{TADS} then adds a
trajectory-derived structural factor on top of this fixed base
ranking,
\[
s_i^{(t)}
=
\mathcal{R}_i^{(t)}a_i^{(t)}
\left(
1+\lambda\,\widetilde{\mathrm{align}}_i^{(t)}
\right).
\]
Setting $\lambda=0$ recovers the same reward--advantage consensus base,
while $\lambda>0$ tests whether representation-level trajectory
alignment improves dynamic selection. Crucially, the trajectory
anchor is defined independently of the reward and advantage signals.
The stability guarantee therefore concerns the structural component
introduced by \textsc{TADS}, not any particular scalar reward, advantage estimator,
policy, or selection-time consensus rule.
